% !TEX root = ../../main.tex
% =====================================================================
% Chapter: Grand Unification; The Chiral Bar--Cobar Adjunction
%
% A typed conditional assembly theorem for the bar--cobar lane
%
% $\Omega^{\mathrm{ch}} \dashv B^{\mathrm{ch}}$
% on $\BrMon^{\mathrm{comp}}_\infty$
%
% whose unit/counit assemble Theorems A, B, C, D, H only after the
% signed log-FM, firewall, completion, trace, OCA, and comparison
% packages have been supplied. The three avatars
%
% L1 = chain-level ensemble (genus-$g$ trace on $U_\cA$)
% L2 = non-Tannakian categorical bar-cobar
% L3 = $(\infty,1)$-adjunction on braided conilpotent factorisation
% coalgebras
%
% are three presentations on their common comparison domain.
%
% Conventions:
% - $U_\cA := \cA \bowtie \cA^!_\infty$ conditional envelope
% (Drinfeld-double-type factorisation bialgebra: bowtie denotes
% factorisation-matched twist, not set-theoretic join)
% - $\BrMon^{\mathrm{comp}}_\infty$ = presentable $\infty$-category
% of conilpotent braided monoidal compact-generated presheaves on
% $\overline{\Conf}(X)$; equivalently, factorisation $\infty$-
% coalgebras on the Ran-compactified configuration space (Lurie
% HA~5.2; Francis-Gaitsgory chiral Koszul; Costello--Gwilliam II)
% - $\Phi_d^{(\Sigma_{d-1},C)}$ = two-stage CY$_d$-to-chiral specialization (Vol III: ch:cy-to-chiral);
% Rozansky-Witten category $\RW_4$ is the genuine output of
% $\Phi_4|_{\mathrm{HK}}$ on hyperkaehler targets.
% - Native-$E_n$ vs derived-$E_n$: native-$E_n$ is the structure
% inherited from factorisation; derived-$E_n$ is the structure
% exhibited by the derived centre. These never coincide on
% non-formal inputs; the Grand Theorem tracks both.
% =====================================================================

\providecommand{\BrMon}{\mathsf{BrMon}}
\providecommand{\BrMoncomp}{\BrMon^{\mathrm{comp}}_{\infty}}
\providecommand{\Coalg}{\mathsf{Coalg}}
\providecommand{\FactCoalg}{\Coalg^{\mathrm{fact}}_{\mathrm{conil}}}
\providecommand{\FactAlg}{\mathsf{ChirAlg}}
\providecommand{\BCh}{B^{\mathrm{ch}}}
\providecommand{\OmCh}{\Omega^{\mathrm{ch}}}
\providecommand{\Uch}{U^{\mathrm{ch}}}
\providecommand{\RW}{\mathsf{RW}}
\providecommand{\HK}{\mathrm{HK}}
\providecommand{\trgenus}{\mathrm{tr}_{g}}
\providecommand{\unitcounit}{\eta \dashv \varepsilon}
\providecommand{\cZch}{\cZ^{\mathrm{der}}_{\mathrm{ch}}}
\providecommand{\barBch}{\bar B^{\mathrm{ch}}}
\providecommand{\Confbar}{\overline{\Conf}}
\providecommand{\ChirHoch}{\mathrm{ChirHoch}}
\providecommand{\kappader}{\kappa}
\providecommand{\Phich}{\Phi^{\mathrm{ch}}}
\providecommand{\Psich}{\Psi^{\mathrm{ch}}}
\providecommand{\Vrep}{\mathsf{Rep}}
\providecommand{\locsys}{\mathrm{LocSys}}
\providecommand{\DMod}{\mathcal{D}}
\providecommand{\BarOrd}{B^{\mathrm{ord}}}
\providecommand{\ChapterIntro}[1]{\par\noindent\textit{#1}\par\medskip}

\chapter{Grand Unification: One Adjunction, Five Theorems, Three Avatars}
\label{chap:grand-unification-platonic}
\index{grand unification!bar-cobar adjunction|textbf}
\index{chiral bar-cobar!universal property|textbf}
\index{braided conilpotent factorisation coalgebras|textbf}
\index{L1 avatar (chain-level ensemble)|textbf}
\index{L2 avatar (non-Tannakian categorical bar-cobar)|textbf}
\index{L3 avatar ($\infty$-categorical adjunction)|textbf}

\ChapterIntro{Theorems A, B, C, D, H and the three avatars L1, L2,
L3 meet only after the source and target of each map have been fixed.
The ordered bar coalgebra, the Verdier--Koszul dual, the derived
chiral centre, the physical bulk, and the scalar modular shadow are
different objects. The adjunction
\[
 \OmCh \;\dashv\; \BCh
 \colon
 \FactAlg_X
 \rightleftarrows
 \BrMoncomp_X
\]
organises the bar--cobar lane. It does not by itself construct the
signed logarithmic Fulton--MacPherson cooperad, the completed
Positselski tower, the open--closed comparison, or the
Calabi--Yau-to-chiral comparison. This chapter records the exact
hypothesis packages under which the separate theorem surfaces assemble
into one typed diagram.}

\section{The Grand Theorem}
\label{sec:grand-theorem}

\subsection{Statement}
\label{ssec:grand-theorem-statement}

\begin{theorem}[Typed assembly theorem for chiral bar--cobar]
\label{thm:grand-unification-platonic}
\ClaimStatusConditional
\textup{[Type: Open, Cat\&Chain, levels $1$--$5$, hypothesis package
$\mathrm{LF}+\mathrm{FW}+H_A+\mathrm{CP}+H_H+H_D+\mathrm{OCA}+H_{L2}$.]}
Fix a smooth projective curve $X/\CC$ and an augmented,
conformal-weight filtered $E_1$-chiral algebra $\cA$ on $X$ with
finite-dimensional weight windows. Assume the following packages.
\begin{enumerate}[label=\textup{(\alph*)}]
\item \emph{Signed log-FM package $\mathrm{LF}$.} The relative
logarithmic Fulton--MacPherson chains form a modular cooperad
$C^{\log\mathrm{FM}}_{\mathrm{mod}}(X)$ with cocompositions
\[
 \Delta_\Gamma^{\log}=(\nu_\Gamma)_*\circ i_\Gamma^!
\]
including determinant-line signs, automorphism normalisations, Gysin
residue maps, proper pushforwards, and homotopy-coherent
coassociativity. With these signs the modular differential
$D_{\mathrm{mod}}^{\log}$ satisfies
$(D_{\mathrm{mod}}^{\log})^2=0$.
\item \emph{Firewall package $\mathrm{FW}$.} Every comparison uses the
five-object firewall of Convention~\ref{conv:master-five-object-firewall}
and Theorem~\ref{thm:typed-verdier-koszul-firewall}:
\[
 \cA,\quad B^{\ord}(\cA),\quad
 \cA^{\mathrm i}=C_X(s^{-1}V,s^{-2}R),\quad
 \cA^!,\quad Z^{\der}_{\ch}(\cA).
\]
The quadratic presentation carries the comparison
$q_\cA\colon\cA^{\mathrm i}\to B^{\ord}(\cA)$; its
quasi-isomorphism is the conclusion of quadratic Koszul recognition.
The counit $\OmCh B^{\ord}(\cA)\to\cA$ is bar--cobar inversion;
Verdier duality produces $\cA^!$; Hochschild cochains produce
$Z^{\der}_{\ch}(\cA)$. These arrows are not interchangeable.
\item \emph{Theorem A package $H_A$.} The hypotheses of
Theorem~\ref{thm:koszul-reflection} hold: augmentation,
augmentation-ideal completion, finite-dimensional conformal-weight bar
pieces, and the weight-completed/pro-conilpotent ambient for
class~$\mathsf M$.
\item \emph{Completed Positselski package $\mathrm{CP}$.} The finite
weight windows of the chiral bar coalgebra form a strict
Mittag--Leffler tower satisfying \textup{(CP1)--(CP3)} of
Theorem~\ref{thm:chiral-positselski-weight-completed}. In particular
all Milnor $\varprojlim^1$ terms in the co/contra comparison are
controlled.
\item \emph{Theorem H package $H_H(\cA;S)$.} A complete chart model
$Q_\cA$ is compared quasi-isomorphically with
$R\!\operatorname{Hom}_{\cA^e}(\cA,\cA)$ and carries a strong
deformation retract onto a complex $K_{\cA,S}$ supported in the
finite set~$S$.  A finite-window construction includes the
incidence-compatible collision retracts, derived inverse-limit
comparison, and ordered-to-symmetric averaging quasi-isomorphism.
\item \emph{Trace package $H_D$.} The modular trace is defined on the
signed log-FM cooperad and is compatible with clutching.  Its four
components are the genus-one trace~$H_D^1$, the virtual Hodge
comparison~$H_D^K$, the derived comparison~$H_D^{\mathrm{tr}}$, and
the stable-graph functional~$H_D^{\mathrm{graph}}$.
\item \emph{Comparison packages.} The physical-bulk reading requires
the open--closed algebra comparison
$\mathrm{OCA}$:
$\mathcal O^{\mathrm{phys}}_{\mathrm{bulk}}(\cA)
\xrightarrow{\sim} Z^{\der}_{\ch}(\cA)$. The categorical and
Calabi--Yau faces require the two-stage functor
$\Phi_d^{(\Sigma_{d-1},C)}
=\SpCh_{\Sigma_{d-1},C}\circ\PhiFA_d$ with its gerbe, Hall,
completion, and descent hypotheses; call this package $H_{L2}$.
\end{enumerate}

\medskip\noindent
\textbf{Two-stage CY-to-chiral notation gate
\textup{(}conditional\textup{)}.}
\label{prop:gu-two-stage-cy-to-chiral-gate}
\index{CY-to-chiral functor!two-stage gate}
\index{Hall--Drinfeld double!notation gate}
\textup{[}Type signature:
$(\mathrm{CY}\to\mathrm{Open},\mathrm{comparison\ package},
\mathrm{level}=0\to 3\to 5,
\mathrm{hypothesis\ package}=H_{L2};
\mathrm{reference\ curve}\ C;
\mathrm{transverse\ manifold}\ \Sigma_{d-1};
\mathrm{factorisation\ homology};
\mathrm{gerbe/Hall/descent/completion})$\textup{]}
The symbol $\Phi_d$ is not used as a one-stage functor on any
load-bearing comparison surface. The typed notation is
\begin{equation}
\label{eq:gu-two-stage-cy-to-chiral}
\Phi_d^{(\Sigma_{d-1},C)}
:=
\SpCh_{\Sigma_{d-1},C}\circ\PhiFA_d .
\end{equation}
Here $\PhiFA_d$ is the first-stage factorisation-algebra output of
the CY$_d$ input; $\SpCh_{\Sigma_{d-1},C}$ is the chiral
specialisation along an oriented transverse manifold
$\Sigma_{d-1}$ followed by restriction to the smooth reference
curve~$C$. The package $H_{L2}$ includes the factorisation-homology
existence, descent, gerbe, completion, and finite-window hypotheses
needed for the composition in \eqref{eq:gu-two-stage-cy-to-chiral}.

The $d=2$ K3 slice and the $d\ge3$ Hall/CoHA slices are separate:
\begin{equation}
\label{eq:gu-cy-slice-separation}
d=2:\ \text{Mukai-Heisenberg/K3 surface lane},\qquad
d\ge3:\ \text{Hall--Drinfeld or CoHA comparison lane}.
\end{equation}
If a shifted Hochschild grading uses an $HH^{-1}$ slice, the convention
is
\begin{equation}
\label{eq:gu-hh-minus-one-zero-slice}
HH^{-1}=0
\end{equation}
unless the shifted complex defining the slice is supplied explicitly.

The Hilbert-scheme Fock-space bridge is a character or state-space
comparison. It is not a categorical lift unless a theorem supplies the
functor on categories and proves preservation of the relevant bar,
centre, or trace structures. The CoHA source is the cohomological Hall
algebra attached to the stated CY/moduli problem in finite charge and
weight windows. The Hall--Drinfeld double is the Drinfeld double of
that windowed CoHA with its product, coproduct, Hopf pairing, and
completion. The double structure is part of $H_{L2}$: it is available
only after product--coproduct compatibility and nondegeneracy of the
pairing have been proved on the same finite windows. The Lusztig
specialisation, when used, is the integral-form root-of-unity
specialisation, including the local convention $\hbar^2=-1/8$ on the
K3/B-family lane; it must commute with the finite-window transition
maps before any completed statement is made.

Then the bar--cobar lane carries the presentable
$\infty$-categorical adjoint pair
\begin{equation}
\label{eq:grand-adjunction-platonic}
 \OmCh \;\dashv\; \BCh
 \colon\;
 \FactAlg_X
 \;\rightleftarrows\;
 \BrMoncomp_X,
\end{equation}
with left adjoint the chiral cobar construction and right adjoint the
chiral bar construction, and the following typed readouts are valid.
\begin{enumerate}[label=\textup{(\roman*)}]
\item \emph{Ordered twisting morphism.} The unit and counit of
\eqref{eq:grand-adjunction-platonic} are represented, on the ordered
chain presentation, by the universal Maurer--Cartan element
\[
 \Theta_\cA\in
 \mathrm{Conv}\bigl(B^{\ord}(\cA),\cA\bigr),
 \qquad
 D\Theta_\cA+\tfrac12[\Theta_\cA,\Theta_\cA]=0,
\]
with $D^2=0$ on the modular logarithmic enhancement only under
$\mathrm{LF}$.
\item \emph{Koszul duality locus.} The full $\infty$-subcategory
$\FactAlg^{\mathrm{Kos}}_X \subset \FactAlg_X$ on which $\varepsilon$
is an equivalence is presented by chirally Koszul algebras and
carries a symmetric duality $\cA \mapsto \cA^!$ induced from
Verdier dualisation of factorisation $\DMod$-modules;
equivalently $\BCh$ restricts there to an equivalence of
presentable stable $\infty$-categories under $H_A$.
\item \emph{Fixed-coalgebra second-kind comparison.} For a supplied
strict completed chiral CDG-coalgebra
$\widehat C=\varprojlim_N C_N$, the package $\mathrm{CP}$ gives
\[
D^{\mathrm{co}}_{\mathrm{ch}}(\widehat C\text{-}\mathrm{comod})
\simeq
D^{\mathrm{ctr}}_{\mathrm{ch}}(\widehat C\text{-}\mathrm{contra}).
\]
This equivalence compares module categories over the same coalgebra.
The quadratic reconstruction map $q_A\colon A^{\mathrm i}\to B_X(A)$
and the raw counit cone have their own type signatures.
\item \emph{Centre-local-system complementarity.} The C0 comparison
identifies the degree-zero strict-flat fibre with the ordinary centre
local system~$\mathcal Z(\cA)$.  The represented involution~$\sigma$
then splits
\[
 \mathbf C_g(\cA)
 =R\Gamma(\overline{\mathcal M}_g,\mathcal Z(\cA))
\]
into its C1 eigensummands.  A supplied brace comparison
$\iota_Z^{\mathrm{der}}$, followed by~$H^0$ and C0, relates this
coefficient system to the Theorem~H derived centre
$Z^{\der}_{\ch}(\cA)$.  The normalized Theorem~D trace gives
$\kappa+\kappa^!\in\{0,13,250/3\}$ on the represented Vol~I
witnesses.  The Bershadsky--Polyakov and Mukai rows carry their
separate comparison packages.
\item \emph{Deformation, Hodge, and graph traces.}  For $g\geq2$, the
native class is
$\operatorname{Obs}^{\mathrm{def}}_g(\cA)\in H^2(\Def_g(\cA))$.
Under~$H_D^1$, the pointed class
$\operatorname{Obs}^{\mathrm{def}}_{1,1}(\cA)$ has trace
$\kappa(\cA)\lambda_1$.  Under~$H_D^K$, the scalar virtual object and
its Hodge image are
\[
 \mathfrak O_g^K(\cA)=\kappa(\cA)\lambda_{-1}(\mathbb E_g),
 \qquad
 \operatorname{ch}_g\!\left(\mathfrak O_g^K(\cA)\right)
 =(-1)^g\kappa(\cA)\lambda_g.
\]
The comparison with the deformation class is~$H_D^{\mathrm{tr}}$.
Under~$H_D^{\mathrm{graph}}$, the genus-$g$ cohomological trace and
graph functional are
\begin{equation}
\label{eq:genus-g-trace-platonic}
 \trgenus^{\mathrm{cl}} \colon
 \BCh(\cA)
 \;\longrightarrow\;
 H^\ast(\overline{\cM}_{g,n},\CC),
 \qquad
 \operatorname{Tr}^{\mathrm{graph}}_g\colon
 H^\ast(\overline{\cM}_{g,n},\CC)\longrightarrow\CC .
\end{equation}
Their composite has numerical value
\[
F_g(\cA)
=F_g^{\mathrm{sc}}(\cA)+\delta F_g^{\mathrm{cross}}(\cA),
\qquad
F_g^{\mathrm{sc}}(\cA)=\kappa(\cA)\lambda_g^{\mathrm{FP}},
\]
on the multi-weight scalar lane.
\item \emph{Family-indexed Hochschild support.} The chiral Hochschild object
$\ChirHoch^\bullet(\cA) := \cZch(\cA)$ is naturally equivalent, inside
the adjunction~\eqref{eq:grand-adjunction-platonic}, to the
endomorphism object
$\End_{\BrMoncomp_X}(\BCh(\cA))$ once the centre-comparison map is
part of the hypothesis package.  Theorem~H assigns a support set
$S$ through $H_{\mathrm H}(\cA;S)$.  Its basic benchmarks are the
Virasoro support $\{0,2,3\}$ and the rank-one even free-superboson
dimensions $(2,1,0,\ldots)$.
\item \emph{Three-avatar presentation.} The adjoint
pair~\eqref{eq:grand-adjunction-platonic} admits three canonical
presentations on the common domain where their comparison packages
hold:
\begin{itemize}
\item \emph{L1 (chain-level ensemble).} Presented on chain-complexes
 in the ordered-configuration model: $\BCh(\cA)$ is $\BarOrd(\cA) =
 T^c(s^{-1}\bar\cA)$ equipped with the Beilinson-Drinfeld
 factorisation differential; $\OmCh$ is the ordered cobar
 construction; unit/counit are the explicit Mittag-Leffler
 witnesses of Chapter~\ref{chap:bar-cobar} and
 Theorem~\ref{thm:chiral-positselski-weight-completed}.
\item \emph{L2 (non-Tannakian categorical).} Presented on the
 presentable stable $\infty$-category
 $\Vrep(\cA) := \cA\text{-}\mathsf{mod}^{\mathrm{fact}}$ as the
 bar--cobar reconstruction on braided module categories; the
 formal-locus restriction recovers the Tannakian reconstruction,
 while the non-formal complement is detected by $\RW_d$-type
 category-valued outputs of the Calabi--Yau specialization
 $\Phi_d^{(\Sigma_{d-1},C)}$ under $H_{L2}$.  Formality of the
 chiral $A_\infty$-structure and quadratic Koszul recognition are
 separate coordinates; the latter is carried by
 $q_A\colon A^{\mathrm i}\to B_X(A)$ for each archetype row.
\item \emph{L3 ($\infty$-categorical, genus-$g$).} Presented on the
 Lurie-Francis-Gaitsgory machinery \cite[HA~5.2]{LurieDAGX}, \cite{FG12}
 of factorisation $\infty$-operads with values in presentable
 stable $\infty$-categories; the genus-$g$ trace is the homotopy
 fixed-point spectrum of the $S^1$-action on $\Uch_\cA :=
 \cA \bowtie_{\infty} \cA^!_\infty$ restricted to
 $\overline{\cM}_{g,n}$ only after the universal-envelope and trace
 packages have been supplied.
\end{itemize}
Each clause retains its displayed source, target, and comparison map.
\end{enumerate}
\end{theorem}

\begin{remark}[Open construction problems]
\label{rem:gu-missing-packages}
\ClaimStatusConditional
The theorem is an assembly theorem.  Its unrestricted form is
equivalent to the following construction problems:
\begin{enumerate}[label=\textup{(\arabic*)}]
\item construct the signed relative log-FM modular cooperad, including
orientation lines, Gysin residues, proper pushforwards, automorphism
normalisations, and homotopy-coherent coassociativity;
\item prove \textup{(CP1)--(CP3)} for every class-$\mathsf M$ tower
used downstream, including the Hom-cone Mittag--Leffler calculation;
\item identify the Hochschild defect complex
$K^\bullet_{\mathrm{DH}}(\cA)$ off the PBW/generic locus and prove
acyclicity before claiming concentration;
\item construct OCA whenever the derived chiral centre is read as the
physical bulk;
\item construct the Hall--Borcherds chain datum
$(\beta,\delta,\varepsilon,\tau)$ before promoting $\Delta_5$ from a
scalar denominator to an operator-level theorem;
\item prove the two-stage CY-to-chiral functor
$\SpCh_{\Sigma_{d-1},C}\circ\PhiFA_d$ with completion, gerbe, and
descent data before using L2 as a theorem rather than a recognition
target.
\end{enumerate}
\end{remark}

\subsection{The five theorem corollaries}
\label{ssec:grand-theorem-five-corollaries}

\begin{corollary}[Theorem~A is the adjunction itself]
\label{cor:gu-theorem-a-platonic}
\ClaimStatusConditional
Under the Theorem A package $H_A$, the bar--cobar clause of
Theorem~\ref{thm:grand-unification-platonic} is exactly
Theorem~\ref{thm:koszul-reflection}: the ordered chiral bar functor
is right adjoint to the ordered chiral cobar functor, and the counit
\[
 \OmCh_X B^{\ord}_X(\cA)\longrightarrow \cA
\]
is a quasi-isomorphism on the Koszul/complete locus.
\end{corollary}

\begin{proof}
Theorem~\ref{thm:koszul-reflection} supplies the Hom-bijection, the
universal twisting morphism, and the counit quasi-isomorphism under
augmentation, completion, and finite-window hypotheses. The assembly
theorem imports exactly those hypotheses as $H_A$ and adds no raw
class-$\mathsf M$ direct-sum statement.
\end{proof}

\begin{corollary}[Theorem~B is quadratic recognition]
\label{cor:gu-theorem-b-platonic}
\ClaimStatusConditional
For a connected positive quadratic chiral presentation, Theorem~B is
carried by the comparison
\[
q_A\colon A^{\mathrm i}\longrightarrow B_X(A).
\]
Under $H_{\mathrm{CL}}(A,A^{\mathrm i},\tau_{\mathrm i})$, its
quasi-isomorphism is equivalent to acyclicity of the two quadratic
twisted tensor products and to concentration of the PBW/bar spectral
sequence on the Koszul diagonal.  The package \textup{(CP1)--(CP3)}
governs the adjacent fixed-coalgebra co--contra comparison.
\end{corollary}

\begin{proof}
This is Theorem~\ref{thm:theorem-B-scope-quadratic-recognition}.
Theorem~\ref{thm:chiral-positselski-weight-completed} has source and
target the second-kind module categories over a fixed completed
coalgebra, and therefore supplies a neighbouring comparison.
\end{proof}

\begin{corollary}[Theorem~C and its normalized scalar trace]
\label{cor:gu-theorem-c-platonic}
\ClaimStatusConditional
Property~\textup{(iv)} of
Theorem~\ref{thm:grand-unification-platonic} is the C0/C1
centre-local-system decomposition.  After the normalized Theorem~D
trace, the computed Vol~I landmark values are
$\{0,13,250/3\}$.  The Mukai arithmetic value~$8$ has its conjectural
chiral interpretation under~$H_{\mathsf B}$, while the
Bershadsky--Polyakov value~$25/3$ remains conditional on its genus-one
curvature calculation.  These statuses are recorded in
Convention~\ref{conv:theorem-c-bucket}.
\end{corollary}

\begin{proof}
Chapter~\ref{ch:landscape-census} and
Chapter~\ref{ch:theorem-C-refinements-platonic} compute the listed
scalar values on their family witnesses.  Theorem~C supplies the C1
eigensummands of~$\mathbf C_g(\cA)$, and Theorem~D supplies their
normalized trace.  The object firewall keeps
$B^{\ord}(\cA)$, $Z^{\der}_{\ch}(\cA)$,
$\mathcal Z(\cA)$, and the resulting scalar in their native targets.
\end{proof}

\begin{corollary}[Theorem~D separates deformation, Hodge, and graph traces]
\label{cor:gu-theorem-d-platonic}
\ClaimStatusConditional
Property~\textup{(v)} of
Theorem~\ref{thm:grand-unification-platonic} gives
$\operatorname{Obs}^{\mathrm{def}}_g\in H^2(\Def_g)$ for $g\geq2$,
$\operatorname{tr}_1\operatorname{Obs}^{\mathrm{def}}_{1,1}
=\kappa\lambda_1$ under~$H_D^1$,
$\mathfrak O_g^K=\kappa\lambda_{-1}(\mathbb E_g)$ and
$\operatorname{ch}_g(\mathfrak O_g^K)=(-1)^g\kappa\lambda_g$ under
$H_D^K$, and, under~$H_D^{\mathrm{graph}}$,
\[
F_g(\cA)=F_g^{\mathrm{sc}}(\cA)
       +\delta F_g^{\mathrm{cross}}(\cA),
\qquad
F_g^{\mathrm{sc}}(\cA)=\kappa(\cA)\lambda_g^{\mathrm{FP}}.
\]
\end{corollary}

\begin{proof}
Theorem~\ref{thm:modular-characteristic} gives these four clauses.
The assembly theorem imports their trace and clutching hypotheses as
$H_D$.
\end{proof}

\begin{corollary}[Theorem~H is family-indexed Hochschild support]
\label{cor:gu-theorem-h-platonic}
\ClaimStatusConditional
For every package $H_H(\cA;S)$, property~\textup{(vi)} of
Theorem~\ref{thm:grand-unification-platonic} identifies
$\ChirHoch^\bullet(\cA)$ with the cohomology of $K_{\cA,S}$ and
places its support in~$S$:
\[
 \operatorname{Supp}\ChirHoch^\bullet(\cA)\subseteq S.
\]
\end{corollary}

\begin{proof}
Property~(vi) names the centre-comparison map from
$\End_{\BrMoncomp_X}(\BCh(\cA))$ to the chiral chart complex.
The strong deformation retract in $H_H(\cA;S)$ computes its
cohomology and support.  The bounded Virasoro support
$\{0,2,3\}$ and rank-one even superboson dimensions
$(2,1,0,\ldots)$ enter through their respective bounded-to-chart
quasi-isomorphisms.
\end{proof}

\subsection{The three-avatar corollaries (L1, L2, L3)}
\label{ssec:grand-theorem-three-avatars}

\begin{corollary}[L1; chain-level ensemble via genus-$g$ trace]
\label{cor:gu-L1-platonic}
\ClaimStatusConditional
Under $\mathrm{LF}$, $H_A$, $\mathrm{CP}$, and $H_D$, the L1 avatar is
the ordered chain presentation of the bar--cobar lane, evaluated by
the modular trace along $\overline{\cM}_{g,n}\to\mathrm{pt}$. On that
scope the displayed identity is
\begin{equation}
\label{eq:L1-avatar-platonic}
 \trgenus \!\left( \OmCh(\BCh(\cA)) \right)
 \;=\;
 \int_{\overline{\cM}_{g,n}} \Theta_{\cA}^{\wedge g}
 \;=\;
 F_g(\cA).
\end{equation}
\end{corollary}

\begin{proof}
The ordered chain model gives $\BCh(\cA)=B^{\ord}(\cA)$ and
$H_A+\mathrm{CP}$ supplies the completed counit comparison. The trace
package $H_D$ is exactly the extra datum required to integrate
$\Theta_\cA^{\wedge g}$ over the modular boundary. Without $H_D$ the
left side remains a bar--cobar object and no genus-$g$ scalar
identity follows.
\end{proof}

\begin{corollary}[L2; non-Tannakian categorical bar--cobar]
\label{cor:gu-L2-platonic}
\ClaimStatusConditional
On the $H_{L2}$ comparison domain and on the non-formal complement,
the $\infty$-adjunction restricts to
a categorical bar--cobar adjunction
\[
 \OmCh^{\mathrm{cat}} \;\dashv\; \BCh^{\mathrm{cat}}
 \colon
 \mathsf{BrCat}^{\mathrm{cat}}_X
 \;\rightleftarrows\;
 \mathsf{BrCat}^{\mathrm{cat,cocomp}}_X
\]
on braided monoidal $\infty$-categories of factorisation modules. In
this domain the formal locus maps to the Tannakian fixed locus. The
non-Tannakian stratum is detected by failure of the categorical bar to
be fibrewise semisimple; the Rozansky-Witten category $\RW_4$ is the
genuine output of $\Phi_4|_{\HK}$ in the constructed hyperkaehler
cases, with non-Tannakian character witnessed by the categorical
structure of $\RW_4$.  The non-Tannakian comparison tracks formality
of the chiral $A_\infty$-structure.  Membership in the Koszul locus is
the family-specific quasi-isomorphism status of~$q_A$.
\end{corollary}

\begin{proof}[Proof sketch and scope]
The categorical bar--cobar of Francis-Gaitsgory
$(\infty,1)$-operadic machinery \cite{FG12} restricts to
$\BrMoncomp_X$ by the presentability of
$\cA\text{-}\mathsf{mod}^{\mathrm{fact}}$; the non-Tannakian
stratum is detected, on the $H_{L2}$ comparison domain, by the
complement of the locus where fibrewise exact $\otimes$-functors into
$\Vrep_\infty$ are representable. On the hyperkaehler targets for
which the Vol~III chiral-to-CY functor $\Phi_4$ has been constructed,
the Rozansky-Witten category $\RW_4$ gives the concrete witness; the
absence of an exceptional collection is supporting evidence, not an
if-and-only-if criterion.
\end{proof}

\begin{corollary}[L3; $(\infty,1)$-categorical adjunction]
\label{cor:gu-L3-platonic}
\ClaimStatusConditional
The Lurie-Francis-Gaitsgory
$(\infty,1)$-presentation of bar--cobar
\cite[HA~5.2.3]{LurieDAGX}, \cite{FG12} is property~(vii.L3) of
Theorem~\ref{thm:grand-unification-platonic}; the two presentations
L1 and L3 are related by factorisation Dold-Kan
\cite[HA~1.2.3]{Lurie17} on their common completed finite-window
domain, and L2 sits between them only under the comparison package
$H_{L2}$.
\end{corollary}

\section{Proof of the Grand Theorem}
\label{sec:grand-theorem-proof}

\subsection{What is assembled}
\label{ssec:grand-theorem-new-content}

The theorem does not assert that a single presentability argument
constructs every object in the programme. It assembles already named
lanes after their hypotheses have been supplied. Clause~(i) is the
ordered Maurer--Cartan bar lane. Clause~(ii) is the Koszul locus of
Theorem~\ref{thm:koszul-reflection}. Clause~(iii) is the completed
Positselski theorem of
Theorem~\ref{thm:chiral-positselski-weight-completed}. Clause~(iv) is
the Theorem~C decomposition of the strict-flat centre-local-system
complex; clause~(v) applies Theorem~D on the trace lane and includes its
scalar shadow. Clause~(vi) is Theorem~H under its PBW
and completion hypotheses. Clause~(vii) records the L1/L2/L3
comparison only on the common domain where the comparison packages
exist.

\begin{proof}[Proof of Theorem~\ref{thm:grand-unification-platonic}]
\emph{Step 1: the bar--cobar adjunction.}
Theorem~\ref{thm:koszul-reflection} and the Francis--Gaitsgory
universal twisting property \cite[\S6]{FG12} give the adjunction
\eqref{eq:grand-adjunction-platonic} in the ordered, augmented,
complete chiral ambient. The finite-window and completion hypotheses
in $H_A$ are exactly the hypotheses under which the counit is a
quasi-isomorphism.

\emph{Step 2: modular logarithmic enhancement.}
The ordered bar differential satisfies the Maurer--Cartan equation on
the fixed-curve Ran model. Passing to modular logarithmic families
adds the planted-forest, collision, sewing, and trace terms. The
assumption $\mathrm{LF}$ supplies the missing determinant-line signs,
Gysin residues, proper pushforwards, automorphism normalisations, and
homotopy-coherent coassociativity; hence
$(D_{\mathrm{mod}}^{\log})^2=0$ is available exactly on that surface.

\emph{Step 3: completed inversion.}
Theorem~\ref{thm:chiral-positselski-weight-completed} identifies the
completed coderived comodule category with the completed
contraderived contramodule category after \textup{(CP1)--(CP3)}. This
is the fixed-coalgebra content of clause~(iii).  For class~$\mathsf M$,
the raw counit is governed by the open cone
$\operatorname{Cone}(\varepsilon_A^{\mathrm{raw}})$.

\emph{Step 4: centre, trace, and scalar lanes.}
Convention~\ref{conv:master-five-object-firewall} and
Theorem~\ref{thm:typed-verdier-koszul-firewall} separate bar--cobar
inversion, Verdier duality, and Hochschild centre formation. The
scalar constants of Theorem~C are imported from
Chapter~\ref{ch:landscape-census}; the trace formula of Theorem~D is
imported from Chapter~\ref{ch:higher-genus}; and Theorem~H is imported
from Theorem~\ref{thm:main-koszul-hoch}. Each import is made only with
its corresponding package $H_H$ or $H_D$.

\emph{Step 5: comparison faces.}
OCA converts the derived centre into a physical bulk when such a
statement is made. The L2 face is supplied by $H_{L2}$ and by
Theorem~\ref{thm:L2-chain-level-gerbe-platonic}; without it, the
non-Tannakian and Calabi--Yau faces are recognition targets rather
than consequences of the bar--cobar adjunction.
\qedhere
\end{proof}

\subsection{Primary references for the $\infty$-categorical surface}
\label{ssec:grand-theorem-primary-references}

The proof uses: Lurie \cite[Higher Algebra~5.2, 5.5]{Lurie17} for
presentable $\infty$-adjoint-functor theory and Dold-Kan; Francis-
Gaitsgory \cite{FG12} for the chiral Koszul duality adjunction at
the $(\infty,1)$-categorical level; Costello-Gwilliam \cite{CG17,CG-vol2}
for the factorisation-algebra presentation on the curve~$X$. The
chain-level witnesses are the standard Beilinson-Drinfeld
\cite{BD04} presentation.

\section{Avatar corollaries expanded}
\label{sec:avatar-corollaries-expanded}

\subsection{L1 as a trace identity on the universal envelope}
\label{ssec:L1-trace-identity}

\begin{proposition}[L1 trace identity]
\label{prop:L1-trace-identity-platonic}
\ClaimStatusConditional
\textup{[Type signature: Open quadrant; universal-envelope and
stable-graph trace presentations; levels $4\to5$; universal-envelope,
completion, and $H_D^{\mathrm{graph}}$ packages.]}
Assume the universal-envelope comparison, the graph package
$H_D^{\mathrm{graph}}$, and
the completion packages of Theorem~\ref{thm:grand-unification-platonic}.
Then the constructed genus-$g$ trace of the envelope
$\Uch_\cA := \cA \bowtie \cA^!_\infty$ agrees with the chain-level integration of
$\Theta_\cA$ along $\overline{\cM}_{g,n}$:
\begin{equation}
\label{eq:L1-trace-identity-platonic}
 \trgenus(\Uch_\cA)
 \;=\;
 F_g(\cA)
 \;=\;
 \kappa(\cA)\,\lambda_g^{\mathrm{FP}}
 \;+\;
 \delta F_g^{\mathrm{cross}},
\end{equation}
with $\delta F_g^{\mathrm{cross}}$ the cross-modular correction.  The
scalar-diagonal graph package gives $\delta F_g^{\mathrm{cross}}=0$.
\end{proposition}

\begin{proof}
The universal-envelope comparison supplies the object
$\Uch_\cA$ and the map from its trace to the ordered bar trace. The
graph package $H_D^{\mathrm{graph}}$ supplies integration over the
signed modular log-FM boundary. Theorem~D gives the displayed
diagonal and cross-channel terms.
\end{proof}

\subsection{L2 as the non-Tannakian recognition surface}
\label{ssec:L2-non-tannakian}

On the $H_{L2}$ comparison domain the non-Tannakian stratum is read
through the locus on which $\Phi_d^{(\Sigma_{d-1},C)}$ produces
category-valued outputs (Vol~III, chapter
\texttt{cy\_to\_chiral.tex}). For $d = 4$ in the constructed
hyperkaehler cases, $\Phi_4^{(\Sigma_3,C)}|_{\HK}$ produces $\RW_4$,
the Rozansky-Witten category; failure of an exceptional collection is
a witness for non-Tannakian behaviour, not a classification theorem.

\begin{proposition}[Non-Tannakian projection]
\label{prop:L2-non-tannakian-platonic}
\ClaimStatusConditional
Conditional on Conjecture~$\Pi_3^{\mathrm{ch}}$ of
Vol~III (chain-level CY-A$_3$), the non-Tannakian stratum of L2
is detected by the conditional image of the non-formal CY$_d$ targets
$T$, i.e.\ the targets for which
$\Phi_d^{(\Sigma_{d-1},C)}(T)$ lies in the non-formal categorical
stratum.  Formality and quadratic Koszul recognition are recorded by,
respectively, the transferred higher products and the comparison map
$q_A$.
\end{proposition}

\begin{proof}[Proof sketch]
Under $\Pi_3^{\mathrm{ch}}$, the CY$_3$-to-chiral correspondence
$\Phi_3^{(\Sigma_2,C)}=\SpCh_{\Sigma_2,C}\circ\PhiFA_3$
is defined on chain-level; its output is chirally formal on the
constructed formal comparison locus (classes G and L), and carries
the displayed higher products on the constructed non-formal comparison
locus (classes C and M).  The Koszul status of every output is computed
by its map~$q_A$.  The formality-locus
restriction of L2 recovers the Tannakian stratum on this comparison
domain; the non-formal complement supplies the detected
non-Tannakian stratum. The Rozansky-Witten category is the witness for
$d=4$ on the constructed hyperkaehler targets.
\end{proof}

\subsection{L3 as a presentable $\infty$-category adjunction}
\label{ssec:L3-infinity}

The L3 presentation is the most flexible for stating universal
properties. The adjunction $\OmCh \dashv \BCh$ in L3 form is
characterised by the $\infty$-universal property: $\BCh(\cA)$ is
the classifying object for conilpotent braided factorisation
coalgebras receiving a chiral twisting morphism from $\cA$. This is
the $(\infty,1)$-categorical statement for the bar--cobar lane. The
trace, centre, physical-bulk, and CY comparison faces enter only
through the packages named in
Theorem~\ref{thm:grand-unification-platonic}.

\section{The typed assembly diagram}
\label{sec:single-infty-adjunction}

Theorem~\ref{thm:grand-unification-platonic} fixes the common
bar--cobar spine:
$\infty$-adjunction
\[
 \OmCh \;\dashv\; \BCh
 \colon\;
 \FactAlg_X
 \;\rightleftarrows\;
\BrMoncomp_X,
\]
whose domain is chiral algebras on $X$, whose codomain is
braided conilpotent factorisation coalgebras on $\Confbar(X)$, and
whose unit/counit control the inversion lane. The categorical target
(L2), the chain-level trace (L1), and the genus-$g$ categorical trace
(L3) are comparison faces of this spine only after their listed
hypotheses have been supplied.

\subsection{Three-volume corollary of the Grand Theorem}
\label{ssec:three-volume-corollary}

\begin{corollary}[Three volumes, three faces of the adjunction]
\label{cor:three-volume-platonic}
\ClaimStatusConditional
The three Volumes of the chiral bar--cobar theory develop
complementary faces of the
adjunction~\eqref{eq:grand-adjunction-platonic}, each using both
L1 and L3 presentations and restricting the categorical face L2
to a different geometric locus. This is a conditional comparison
statement: the Volume~I and Volume~II faces below are valid in their
stated scopes, and the Volume~III CY$_3$ gerbe-twisted face has
exactly the conditional strength of
Theorem~\ref{thm:L2-chain-level-gerbe-platonic}. Volume~I develops the chain-level
(L1) and $(\infty,1)$-categorical (L3) presentations on a smooth
projective curve $X$, with the ordered bar complex
$\BarOrd(\cA) = T^c(s^{-1}\bar\cA)$ as the chain-level anchor.
Volume~II develops the L3 presentation on
$\mathrm{SC}^{\mathrm{ch,top}}$-algebras over
$\Confbar(X) \times \bR^k$, with transverse topological structure
extending the L3 machinery. Volume~III develops L2 in the
CY-to-chiral direction via the two-stage specialization
$\Phi_d^{(\Sigma_{d-1},C)}$, and is the source
of the non-Tannakian stratum of L2 as
Theorem~\ref{thm:L2-chain-level-gerbe-platonic} exhibits for $d = 3$.
\end{corollary}

\begin{proof}
The statement is obtained by applying the three comparison corollaries
with their hypotheses left visible. Volume~I supplies the ordered
chain model and the completed bar--cobar adjunction. Volume~II enters
only after the transverse topological $\SCchtop$ comparison is
provided. Volume~III enters only through
$\Phi_d^{(\Sigma_{d-1},C)}$ and the gerbe-twisted L2 package of
Theorem~\ref{thm:L2-chain-level-gerbe-platonic}. Thus the corollary
records a diagram of conditional restrictions between the three
volumes.
\end{proof}

\section{Sources and comparison maps}
\label{sec:grand-theorem-verification}

The numerical inputs imported into the typed assembly have their own
verification sources. The assembly itself remains conditional on the
packages in Theorem~\ref{thm:grand-unification-platonic}.
\begin{itemize}
\item \emph{Source for property~(i).} Chain-level: the
universal twisting morphism $\Theta_\cA$ is constructed explicitly
in Chapter~\ref{chap:bar-cobar}.
The Francis--Gaitsgory reduced-operad equivalence supplies the
categorical ambient; $H_{\mathrm{fact}}$ carries the comparison with
the chiral factorization subcategory.
\item \emph{Source for property~(iii).}  Positselski's theorem gives
the fixed-coalgebra coderived/contraderived equivalence.  Chiral
transport is the separate twisting-data package of
Chapter~\ref{ch:theorem-B-scope-platonic}.
\item \emph{Source for property~(iv).} The computed Vol~I scalar set
is $\{0,13,250/3\}$.  The Bershadsky--Polyakov value is open pending
its genus-one minimal-DS trace, and the Mukai integer~$8$ becomes a
chiral conductor under~$H_{\mathsf B}$; see
Convention~\ref{conv:theorem-c-bucket}.
\item \emph{Source for property~(v).}  The genus-one trace,
virtual Hodge object, signed Hodge image, and stable-graph functional
are the four conditional clauses of
Theorem~\ref{thm:modular-characteristic}.  The genus-one and clutching
chapters construct the corresponding comparison maps.
\item \emph{Source for property~(vi).}  The family package
$H_H(\cA;S)$ supplies the chart retract and hence support in~$S$.
Bakalov--De Sole--Kac supply the bounded Virasoro and rank-one even
free-superboson benchmarks; chart transport is a separate comparison.
\end{itemize}

\section{Open threads}
\label{sec:grand-theorem-open-threads}

The L1 and L3 faces close only on the completed finite-window domain
specified above. The L2 face is a comparison target until the
Costello--Li curved-HKR transfer \cite{CL20}, the gerbe/descent
package, and the two-stage CY-to-chiral specialization have been
constructed for the chosen compact CY$_3$ target. The theorem below
records that conditional chain-level scope; Vol~III
(\texttt{cy\_to\_chiral.tex}) contains the corresponding CY-side
formulation.

\begin{theorem}[L2 chain-level functor into gerbe-twisted braided monoidal $\infty$-category]
\label{thm:L2-chain-level-gerbe-platonic}
\ClaimStatusConditional
Let $X$ be a compact CY$_3$ without a full exceptional collection
(Bondal-Van den Bergh non-representability \cite{BondalVdB03});
let $\alpha_X \in H^3(X,\CC^\times)$ denote the Brauer-type gerbe
class obstructing strictification of the braided $\infty$-category
of perfect complexes to an honest bialgebra; let
$\mathrm{HKR}_\infty^{\mathrm{curv},\hbar}$ denote the
curved-$L_\infty$ quasi-isomorphism of Costello-Li \cite{CL20}.
Conditional on the BCOV transfer of \cite{CL20} extending to the
genus range used on $X$, and on the gerbe/descent data of
\S\ref{ssec:pr-alpha-X-local-global}, there is a chain-level
$\infty$-functor
\begin{equation}
\label{eq:Phi-3-ch-platonic}
 \Phi_3^{\mathrm{ch}}
 \;\colon\;
 \mathrm{Perf}^{\mathrm{CY3}}_\infty(X)
 \;\longrightarrow\;
 \BrMoncomp_{X,\alpha_X}
\end{equation}
valued in the gerbe-twisted presentable $\infty$-category of
conilpotent braided-monoidal compact-generated factorisation
coalgebras on $\Confbar(X)$. The functor~\eqref{eq:Phi-3-ch-platonic}
satisfies three properties.
\begin{enumerate}[label=\textup{(\roman*)}]
\item \emph{Chain-level realisation.} The underlying chain-level
model of $\Phi_3^{\mathrm{ch}}(X)$ is the explicit chain-level bar
$B^{\mathrm{cat}}(D^b(X))$ (\S~\ref{ssec:explicit-bar-model}),
with $L_\infty$-structure inherited from
$\mathrm{HKR}_\infty^{\mathrm{curv},\hbar}$ and with
gerbe-twisting $\alpha_X$ acting on the associator.
\item \emph{Gerbe compatibility.} The class
$\alpha_X \in H^3(X,\CC^\times)$ is the specified obstruction datum
for lifting $\Phi_3^{\mathrm{ch}}$ from the gerbe-twisted target
$\BrMoncomp_{X,\alpha_X}$ to $\BrMoncomp_X$. A strict untwisted lift
exists only after a separate descent computation trivialises this
class.
\item \emph{Categorical-H compatibility.} The ambient Hochschild
complex $\mathrm{ChirHoch}^\bullet(\Phi_3^{\mathrm{ch}}(X))$ is
equivalent to the Costello-Li BCOV ambient Hochschild complex,
concentrated in the three-step filtration of Theorem~H only under the
package $H_H$ of Theorem~\ref{thm:grand-unification-platonic}.
\end{enumerate}
Consequently, on the subcategory of compact CY$_3$ targets for
which the Costello-Li BCOV transfer and descent package are
established, the non-Tannakian stratum of the L2 avatar of
Theorem~\ref{thm:grand-unification-platonic} is detected by
the non-formal chain-level CY$_3$ targets in the conditional image of
$\Phi_3^{\mathrm{ch}}$.
\end{theorem}

\begin{proof}
\emph{Step 1 (lift the derived category to a chain-level $\infty$-functor).}
On a compact CY$_3$, $D^b(X)$ admits no full exceptional
collection \cite{BondalVdB03}, and the derived-to-bialgebra
construction via semiorthogonal decomposition fails. By
\cite[HA~5.1.2.2, 5.1.4.2]{LurieDAGX}, the $\infty$-category
$\mathrm{Perf}^{\mathrm{CY3}}_\infty(X)$ of perfect complexes on a
CY$_3$ is symmetric-monoidally enriched over
$(\Coalg^{\mathrm{fact}}_{\mathrm{conil}})_\infty$; its underlying
braided monoidal structure lands in $\BrMoncomp_X$ up to a
Brauer-type twist.

\emph{Step 2 (identify the gerbe obstruction).}
The Brauer-type twist is controlled by $H^3(X,\CC^\times)$, the
standard obstruction group for strictifying a braided
$\infty$-category to an honest bialgebra \cite{Toen07}, giving the
specified class $\alpha_X = [\mathrm{Gerbe}(\Phi_3)]$. A full
exceptional collection is a sufficient strictification datum in cases
where the semiorthogonal bialgebra construction is available; absence
of such a collection does not by itself compute $\alpha_X$. Thus the
target is $\BrMoncomp_{X,\alpha_X}$ until a separate descent
computation proves $\alpha_X=0$.

\emph{Step 3 (construct the $L_\infty$-structure via BCOV transfer).}
The Costello-Li curved $L_\infty$-quasi-isomorphism
$\mathrm{HKR}_\infty^{\mathrm{curv},\hbar}$ of \cite{CL20}
transfers the Dolbeault-BCOV $L_\infty$-structure on
$\mathrm{PV}^\bullet(X)[[\hbar]]$ to the chain-level
$\infty$-operadic structure on
$\mathrm{Perf}^{\mathrm{CY3}}_\infty(X)$ when the $H_{L2}$ transfer
datum has been supplied. The quintic discussion below records the
finite genus range used in the model calculation
(\S~\ref{ssec:BCOV-transfer-X5}); extension to another compact
CY$_3$ is part of the theorem's scope condition, not a consequence of
the local model alone. Combining with Step~2 gives the chain-level
$L_\infty$-data of
$\Phi_3^{\mathrm{ch}}(X) \in \BrMoncomp_{X,\alpha_X}$.

\emph{Step 4 (verify the object is well-defined).}
The conditional object requires three checks: compact generation of the
chosen chain model, conilpotence of the categorical bar coproduct, and
factorisation compatibility on $\Confbar(X)$. These checks are part of
the $H_{L2}$ package; \S~\ref{ssec:BCOV-transfer-X5} and
\S~\ref{ssec:explicit-bar-model} record the model calculations rather
than an all-CY$_3$ theorem.

\emph{Step 5 (paste to the explicit bar object).}
The $\infty$-functor $\Phi_3^{\mathrm{ch}}$ factors through the
explicit chain model
$B^{\mathrm{cat}}(D^b(X)) \in \BrMoncomp_{X,\alpha_X}$ after the
descent datum of Theorem~\ref{thm:pr-alpha-X-local-global} has been
chosen. The resulting diagram of $\infty$-functors
\begin{equation}
\label{eq:Phi-3-ch-commutative-platonic}
\begin{tikzcd}[column sep=large, row sep=large]
\mathrm{Perf}^{\mathrm{CY3}}_\infty(X)
 \arrow[r, "\Phi_3^{\mathrm{ch}}"]
 \arrow[d, "\mathrm{HKR}_\infty^{\mathrm{curv},\hbar}"']
&
\BrMoncomp_{X,\alpha_X}
 \arrow[d, "\mathrm{pr}_{\alpha_X}"]
\\
\mathrm{PV}^\bullet(X)[[\hbar]]
 \arrow[r, "B^{\mathrm{cat}}"']
&
\BrMoncomp_{X}
\end{tikzcd}
\end{equation}
commutes in the effective descent regime and otherwise commutes only
after the specified Cech-coherent descent homotopy. The upper path is
the L2 realisation of
Theorem~\ref{thm:grand-unification-platonic}; the lower path is the
chain-level realisation via BCOV and the explicit bar. The right
vertical $\mathrm{pr}_{\alpha_X}$ is the local gerbe-trivialising
projection after tensoring with a chosen local trivialisation of
$\alpha_X$. Francis-Gaitsgory universality \cite[Thm~6.3]{FG12}
together with Lurie \cite[HA~5.5.2.9]{LurieDAGX} supplies the
comparison once this descent datum is part of $H_{L2}$; it does not
annihilate the gerbe obstruction.

Property~(i) follows from Step~5. Property~(ii) is Step~2 together
with the descent-scoped projection of
diagram~\eqref{eq:Phi-3-ch-commutative-platonic}. Property~(iii)
identifies $\mathrm{ChirHoch}^\bullet$ with the Costello-Li BCOV
ambient Hochschild complex only on the $H_H+H_{L2}$ comparison
surface, combining Theorem~\ref{thm:theorem-h-on-koszul-locus} with
the BCOV transfer datum of Step~3. On the scope of the theorem, the
non-Tannakian stratum of Corollary~\ref{cor:gu-L2-platonic} is
detected by non-trivial gerbe or non-formal chain data in the
conditional image of $\Phi_3^{\mathrm{ch}}$.
\end{proof}

The two supporting inscriptions referenced in the proof are
recorded in the next subsections.

\subsection{Chain-level BCOV transfer on $X_5$}
\label{ssec:BCOV-transfer-X5}

The curved-$L_\infty$ quasi-isomorphism
$\mathrm{HKR}_\infty^{\mathrm{curv},\hbar}$ of \cite{CL20}
transfers the Dolbeault-BCOV action functional on
$\mathrm{PV}^\bullet(X_5)[[\hbar]]$ to the chain-level
$\infty$-operadic structure on
$\mathrm{Perf}^{\mathrm{CY3}}_\infty(X_5)$. The
Kodaira-Spencer $L_\infty$-algebra
$\mathrm{KS}(X_5) = (\mathrm{PV}^\bullet(X_5), \bar\partial, \{-,-\}_{\mathrm{SN}})$
admits a curved extension at every genus $g \leq g_{\max}(X_5)$
within the Costello-Li gauge; the induced $\infty$-structure on
$\mathrm{Perf}^{\mathrm{CY3}}_\infty(X_5)$ is the data passed to
Step~3 of the proof. Extension to arbitrary compact CY$_3$ is the
scope condition of
Theorem~\ref{thm:L2-chain-level-gerbe-platonic}.

\subsection{Explicit $B^{\mathrm{cat}}(D^b(X_5))$}
\label{ssec:explicit-bar-model}

Given the transfer datum of \S~\ref{ssec:BCOV-transfer-X5}, the bar
construction $B^{\mathrm{cat}}(D^b(X_5))$ is the explicit
chain-level coalgebra whose generators are the perfect complexes on
$X_5$ modulo quasi-isomorphism, whose coproduct is the
Fulton-MacPherson boundary of $\Confbar(X_5)$, and whose compatibility
with the Costello-Li BCOV structure is part of that datum. The
associator is twisted by $\alpha_{X_5}$; after the conilpotence and
factorisation checks included in $H_{L2}$,
$B^{\mathrm{cat}}(D^b(X_5))$ lies in
$\BrMoncomp_{X_5, \alpha_{X_5}}$.

Theorem~\ref{thm:L2-chain-level-gerbe-platonic} scopes the earlier
conditional formulation: the non-Tannakian surface is a chain-level
recognition statement on the subcategory of compact CY$_3$ targets
where the Costello-Li BCOV transfer and gerbe descent data are
supplied, with the gerbe twist retained in the target category.

\subsection{Local-global gluing of the gerbe-trivialising projection}
\label{ssec:pr-alpha-X-local-global}

The projection $\mathrm{pr}_{\alpha_X} : \BrMoncomp_{X,\alpha_X}
\to \BrMoncomp_X$ appearing in
diagram~\eqref{eq:Phi-3-ch-commutative-platonic} is defined
locally: on an open cover $\{U_i\}$ of $X$ over which
$\alpha_X|_{U_i}$ trivialises, tensoring with a local
trivialisation $\tau_i : \alpha_X|_{U_i} \xrightarrow{\sim} 0$
produces a gerbe-trivialised restriction
$\BrMoncomp_{U_i, \alpha_X|_{U_i}} \xrightarrow{\sim}
\BrMoncomp_{U_i}$. On overlaps $U_i \cap U_j$, the two local
trivialisations differ by a $\mathbb{C}^\times$-gerbe cocycle
$g_{ij} \in Z^1(U_i \cap U_j, \underline{\mathbb{C}^\times})$
whose cohomology class $[\{g_{ij}\}] \in \check H^3(X,
\underline{\mathbb{C}^\times})$ is $\alpha_X$ itself. The
question is whether the locally defined functors glue to an
honest $\infty$-functor on $X$, or only up to Čech-coherent
homotopy with obstruction in the Postnikov tower.

\begin{theorem}[Conditional local-global gluing of $\mathrm{pr}_{\alpha_X}$]
\label{thm:pr-alpha-X-local-global}
\ClaimStatusConditional
Let $X$ be a smooth projective variety and
$\alpha_X \in H^3(X, \underline{\mathbb{C}^\times})$ the gerbe
class under consideration. Fix a hypercover $U_\bullet\to X$ on which
$\alpha_X$ is trivialised. The local functors
$\mathrm{pr}_{\alpha_X|_{U_n}}$ glue as follows.
\begin{enumerate}[label=\textup{(\alph*)}]
\item \emph{Effective descent regime.} If a finite \'{e}tale Galois
cover $\pi:\widetilde X\to X$ is given, $\pi^*\alpha_X=0$, and the
resulting descent datum in
$\mathrm{Fun}_\infty(\BrMoncomp_{\widetilde X},
\BrMoncomp_{\widetilde X})$ is effective, then
$\mathrm{pr}_{\alpha_X}$ exists as an honest global
$\infty$-functor
\[
 \BrMoncomp_{X,\alpha_X}\longrightarrow \BrMoncomp_X .
\]
\item \emph{Hypercover regime.} Without such an effective descent
datum, the local projections determine only an object of the
totalisation
\[
 \mathrm{Tot}\,
 \mathrm{Fun}_\infty(\BrMoncomp_{U_\bullet,\alpha_X},
                    \BrMoncomp_{U_\bullet}).
\]
The obstruction to strict globalisation is the class of this
totalisation point in the descent Postnikov tower. It vanishes exactly
when the local projections are equivalent to a global
$\infty$-functor on $X$.
\end{enumerate}
Consequently, diagram~\eqref{eq:Phi-3-ch-commutative-platonic}
commutes strictly only in the effective descent regime. In the
hypercover regime it commutes only after choosing the indicated
descent homotopy, and the obstruction class is part of the L2 data.
\end{theorem}

\begin{proof}
\emph{Regime (a).} On $\widetilde X$ the class $\pi^*\alpha_X$
vanishes, so the twisted and untwisted targets are identified by the
chosen trivialisation. Effectivity of the Galois descent datum is
precisely the hypothesis that this local functor descends along the
\v Cech nerve of $\pi$. Giraud descent for gerbes
\cite[III.3.3.2]{Giraud1971} and Lurie descent for presentable
$\infty$-categories \cite[HA~1.3.1.7]{LurieDAGX} then give the global
functor.

\emph{Regime (b).} For an arbitrary hypercover, the local projections
form a cosimplicial diagram of $\infty$-functors. Joyal--Lurie descent
\cite[HTT~6.5.3.12]{LurieHTT} identifies global functors with the
descent-effective points of its totalisation. The Postnikov tower of
this totalisation records the successive homotopy-coherence
obstructions. Thus strict globalisation is equivalent to vanishing of
that obstruction tower. The commutativity of
diagram~\eqref{eq:Phi-3-ch-commutative-platonic} is therefore a
chosen descent homotopy, not an automatic consequence of the
bar--cobar adjunction.
\end{proof}

\begin{remark}[Strict compact CY$_3$ and the quintic]
\label{rem:pr-alpha-X-quintic}
If the Vol~III CY-side construction supplies, on the quintic
$X_5 \subset \mathbb{P}^4$, a gerbe class
$\alpha_{X_5} = [5 \cdot \mathrm{vol}] \in J^3(X_5) \subset
H^3(X_5, \underline{\mathbb{C}^\times})$ lies entirely in the
continuous Griffiths part
(Proposition~\ref{V3-prop:cy-c-beyond-joyce-three-CY3}(i) of
Vol~III, \texttt{cy\_c\_beyond\_k3e\_existence\_obstruction.tex})
and has infinite order, then Regime~(b) of
Theorem~\ref{thm:pr-alpha-X-local-global} applies:
$\mathrm{pr}_{\alpha_{X_5}}$ exists only up to Čech-coherent
homotopy, and diagram~\eqref{eq:Phi-3-ch-commutative-platonic}
commutes only after the chosen descent homotopy unless the obstruction
tower vanishes. Regime~(a) applies to polarised K3 surfaces only after
one supplies an effective finite descent datum for
$\alpha_{K3}^{\mathrm{tors}} \in
\mathrm{Br}(K3)$ and the Caldararu--Azumaya gerbe; the two regimes
separate the two descent problems in the Joyce--Brauer dictionary
(Remark~\ref{V3-rem:cy-c-beyond-quintic-joyce-not-brauer} of
Vol~III).
\end{remark}

\begin{conjecture}[Universal envelope as ambient]
\label{conj:universal-envelope-ambient-platonic}
\ClaimStatusConjectured
The universal envelope $\Uch_\cA = \cA \bowtie \cA^!_\infty$ is
the final object in the $\infty$-category of factorisation
$\infty$-bialgebras on $X$ which pair $\cA$ with a
factorisation-compatible Koszul-dual companion; equivalently,
every such pairing factors through $\Uch_\cA$.
\end{conjecture}

\begin{conjecture}[Genus-$g$ universality]
\label{conj:genus-g-universality-platonic}
\ClaimStatusConjectured
For every $g \geq 2$ and every class-M chiral algebra $\cA$, the
genus-$g$ trace $\trgenus(\Uch_\cA)$ admits a uniform-weight
decomposition whose cross-modular correction
$\delta F_g^{\mathrm{cross}}$ is a sum of
Fulton-MacPherson boundary contributions; equivalently, on
class~M the shadow-tower depth $r_{\max} = \infty$ forces a
non-trivial FM boundary contribution at every $g \geq 2$.
\end{conjecture}

\section{Reader's concordance: the unification map}
\label{sec:grand-theorem-concordance}

\begin{table}[h]
\centering
\begin{tabular}{l|l|l}
\hline
\textbf{Grand Theorem property} & \textbf{Classical theorem} & \textbf{Location} \\
\hline
(i) Universal twisting morphism & Theorem~A (bar--cobar) & Ch.~\ref{chap:bar-cobar} \\
(ii) Koszul duality locus & Koszul-locus discipline & Ch.~\ref{chap:koszul-pair-structure} \\
(iii) Fixed-coalgebra co--contra comparison & Positselski second-kind equivalence & Ch.~\ref{ch:theorem-B-scope-platonic} \\
(iv) Centre-local-system complementarity and normalized trace & Theorem~C followed by Theorem~D & Ch.~\ref{ch:theorem-C-refinements-platonic} \\
(v) Deformation/Hodge/graph comparison & Theorem~D under $H_D^1,H_D^K,H_D^{\mathrm{tr}},H_D^{\mathrm{graph}}$ & Ch.~\ref{ch:higher-genus} \\
(vi) Hochschild concentration & Theorem~H on the Koszul/generic surface & Thm.~\ref{thm:theorem-h-on-koszul-locus} \\
(vii.L1) Chain-level ensemble & chain-level $\Theta_\cA$ & Ch.~\ref{ch:holographic-datum-master} \\
(vii.L2) Non-Tannakian categorical & non-formal / $\RW_d$-valued & Vol~III, \texttt{cy\_to\_chiral.tex} \\
(vii.L3) Presentable $\infty$-adjunction & Francis-Gaitsgory machinery & \cite[\S6]{FG11}, \cite[HA~5.2]{LurieDAGX} \\
\hline
\end{tabular}
\caption{The unification map: seven properties of the Grand
Theorem, their classical-theorem shadows, and their locations in
the manuscript.}
\label{tab:gu-unification-map-platonic}
\end{table}

The Grand Theorem
\ref{thm:grand-unification-platonic} records the structural arc of
Volume~I as a typed assembly diagram: the bar--cobar adjunction is the
spine, and every scalar, categorical, physical, or Calabi--Yau
projection enters only through its named comparison package.
